\chapterOrPart{Making drawings and plots}\label{ch:drawing}
This chapter contains some examples of figures and charts. More will be added in the future.

It's possible to make all kinds of drawings in \LaTeX\ with the package \verb|TikZ|. The style of objects can be predefined. See for instance the source code for \cref{fig:3systems} where the blue fill colour and the red dashed border of the three systems are defined with the command \verb|\tikzstyle| in the style file \verb|myDrawingStuff.sty|.

\section{Simple TikZ drawing}

\begin{frame}{An example of a simple TikZ drawing}
    \begin{figure}[htbp]\centering
		\begin{tikzpicture}
    		%\draw [help lines] (0,0) grid (10,4);
    		\draw [system] (0,0) ellipse (1.5cm and 1cm) node [black]{open};
    		\draw [system] (4,0) ellipse (1.5cm and 1cm) node [black]{closed};
    		\draw [system] (8,0) ellipse (1.5cm and 1cm) node [black]{isolated};
    		\draw [mass] [<->] (0,0.75) -- + (0,0.5) node [above] {mass};
    		\draw [heat] [<->] (0,-0.75) -- + (0,-0.5) node [below] {energy};
    		\draw [heat] [<->] (4,-0.75) -- + (0,-0.5) node [below] {energy};
		\end{tikzpicture}
    	\caption{Example of some objects with predefined properties defined with the \texttt{tikzstyle} command; here: three kinds of thermodynamic systems}\label{fig:3systems}
    \end{figure}
\end{frame}

\section{Drawing a block diagram}
It's also easy to draw block diagrams as in \cref{fig:control}. The original source code was found on the website \href{www.texample.net/tikz/examples/control-system-principles}{TeXample.net}, built by \cite{texample} and was adapted.

\begin{figure}[ht]\centering
    \begin{tikzpicture}[auto, node distance=2.5cm,>=latex']
        %\draw [helpline] (0,2) grid (10,-3);
        \tikzstyle{block}   = [draw, fill=green!20, rectangle, minimum height=3em, minimum width=6em]
        \tikzstyle{sum}     = [draw, fill=green!20, circle, node distance=1cm]
        \tikzstyle{input}   = [coordinate]
        \tikzstyle{output}  = [coordinate]
        % We start by placing the blocks
        \node [input, name=input] {};
        \node [sum, right of=input] (sum) {};
        \node [block, right of=sum] (controller) {controller};
        \node [block, right of=controller,node distance=4cm] (system) {system};
        % position the node u between the controller and system block
        % we can then position the measurement block below the node u
        \draw [->] (controller) -- node[name=u] {$u$} (system);
        \node [block, below of=u] (measurements) {measurements};
        \node [output, right of=system] (output) {};
    
        % Connecting the nodes 
        \draw [draw,->] (input) -- node[pos=0.2] {$i$} node[pos=0.8] {$+$} (sum);
        \draw [->] (sum) -- node {$e$} (controller);
        \draw [->] (system) -- node [name=y] {$y$}(output);
        \draw [->] (y) |- (measurements);
        \draw [->] (measurements) -| node[pos=0.90] {$-$} 
            node [near end] {$y_\up{m}$} (sum);
        \draw [<-] (system.north) --  +(0,0.4) node[above] {$\delta$};
    \end{tikzpicture}
	\caption{Example of a simple system control diagram drawn in TikZ}\label{fig:control}
\end{figure}

\section{Proper sizing of a figure on a slide}
The aspect ratio of the A4 page format differs from the aspect ratio of a slide in beamer. As a consequence, choosing the proper dimensions of figures, to make them look good on A4 and also on a slide, is sometimes challenging. The trick is to define both the height and width, in combination with the option \verb|keepaspectratio|. It will freeze the aspect ratio of an included picture, to make it fit on A4 and on a slide. Inspect the source code of \cref{fig:radiation} to find out how.

\begin{frame}{The size of this picture depends on the type of document (A4 or slide)}
    \begin{figure}[hb]\centering
		     % specifying the height AND the width of an included picture will stretch the picture.
		     % the keepaspectratio option is needed to avoid this.
		     % on A4, the width=0.7\textwidth is the hardest constraint and it will determine the size
		     % with beamer, the height=0.6\textheight is the hardest constraint and it will determine the size
		     \includegraphics[width=0.7\textwidth,height=0.7\textheight,keepaspectratio]{figs/radiation}
			% \includegraphics[width=0.8\textwidth,height=0.8\textheight,keepaspectratio]{figs/radiation} % remove comment to see the bad result in beamer: picture too high and caption falls off the slide
			% \includegraphics[width=0.8\textwidth,height=0.6\textheight]{figs/radiation} % remove comment to see the stretch
		\caption{This included picture should fit correctly on A4 page and on a slide; here: exchange of thermal radiation for some common geometries, source: \cite{tfs}}\label{fig:radiation}
    \end{figure}
\end{frame}

\section{Adding annotations on an included picture}
An example of an included picture file (\texttt{jpg, png, \dots}) with annotations added in \LaTeX\ is shown in \cref{fig:valve}. The image of the expansion valve is included with the \verb|\includegraphics| command. The dotted CV line, the arrows and symbols are added within \LaTeX\ with simple TikZ drawing commands. Inspect the source file to see how.

\begin{frame}{With \LaTeX\ you can easily add annotations to imported picture files}
    \begin{figure}[htbp]\centering
    	    \begin{tikzpicture}[scale=0.9]\node at (0,0) [above right] {\includegraphics[width=5.04cm]{figs/expansionvalve}};
    			%\draw [help lines] (-2,-1) grid (8,6); % uncomment to show grid
    			%\draw [cv] (2,0) -- (3.5,0) -- ++(0.5,1) -- ++(0,3.5) -- (0,4.5) node[left] {CV}  -- (0,3) -- cycle;
    			\draw [mass][->] (2.75,-0.5) node [below] {$h_1$} to +(0,0.8);
    			\draw [mass][->] (0.2,3.6)  to  +(-0.7,0) node [left] {$h_2$};
    			\draw [heat][->] (2.7,2)  to [out=230, in=0] (0.5,0.5) node [left] {$\qdot\approx 0$};
    			\draw [work][->] (2.7,3) to (0,2) node [left] {$\wdot=0$};
    			\draw [pin] (4.3,2) -- + (1,-2) node [right,align=left] {pressure\\compensation line};
    			\draw [pin] (5.3,3.5) -- + (1,-2) node [right,align=left] {temperature\\sensor bulb};
			\end{tikzpicture}
    	\caption{Annotations made with TikZ on an included picture file; here: a thermostatic expansion valve}\label{fig:valve}
    \end{figure}
\end{frame}

\section{Drawing free-body diagrams}
An example of a mechanical sketch of two bodies on a slope connected with a wire and the related free-body diagrams are shown in \cref{fig:freebody}. The original TikZ-code was written by \cite{freebody} and can be found on \href{http://www.texample.net/tikz/examples/free-body-diagrams}{TeXample net}. Drawings in TikZ can be parametric. In this case, the inclination angle of the slope can be altered by changing the value of \verb|iangle| in the source code. The angle in \cref{fig:freebody} is \ang{45} while the angle in the original drawing was \ang{30}. You can compare both drawings to see the effect.

\begin{frame}{Example of a mechanical drawing}
    \begin{figure}[ht]\centering
        \input{figs/freebody.tex} % source: http://www.texample.net/tikz/examples/free-body-diagrams
        \caption{A mechanical sketch and free-body diagrams made with TikZ, adapted from \cite{freebody}}\label{fig:freebody}
    \end{figure}
\end{frame}

\section{Drawings with a 3D look}
Drawings that create a 3D look, are also easy to make, as shown in \cref{fig:torque}.

\begin{frame}{Combining a drawing with equations next to it}
    \begin{figure}[htbp]\centering
    		\begin{tikzpicture}[scale=0.7]
    % 		\only<article>{\begin{tikzpicture}[scale=1]}
    			% \draw [helpline] (-2,-2) grid (5,5);
    			\draw[fill=blue!20](0,0) circle (1 and 2);
    			\draw[top color=blue!10,bottom color=black!70,middle color=blue!30] (-0.5,2) arc (90:270:1 and 2) -- ++(0.5,0) arc (-90:-270:1 and 2) -- cycle;
    			\draw[top color=white,bottom color=black!70] (0,3mm) arc (90:270:1.5mm and 3mm)--++(2cm,0) arc (-90:-270:1.5mm and 3mm)-- cycle;
    			\draw (2cm,3mm) arc (90:-90:1.5mm and 3mm);
    			\draw[->] (1.5,0.4) arc (55:325:0.2cm and 0.6cm) node[pos=0.8, below] {$\omega$};
    			\draw[measure,->] (0,0) --  ++(155:1.4) node[midway, above] {$\vec{r}$};
    			\draw[force,->] (0,0) ++(155:1.4) --  ++(255:2.5) node[near end, left] {$\vec{F}$};
    			\draw[torque,->] (0,0) -- ++(3,0) node[right,align=left] {$\vec{T}=\vec{r}\times \vec{F}$};
    			\node at (6,1.5) [right] {$W_\up{shaft}=T\,\Delta\theta\units{J}$};
    			\node at (6,0.5) [right] {$\wdot_\up{shaft}= T\,\omega=T\, \frac{2\pi N}{60}\units{W}$};
    			\node at (6,-0.5) [right] {$T$};
    			\node at (6,-1.5) [right]   {$\Delta\theta$};
    			\node at (6,-2.5) [right] {$N$};
    			\node at (7,-0.5) [right] {torque [Nm]};
    			\node at (7,-1.5) [right]   {rotation angle [rad]};
    			\node at (7,-2.5) [right] {rotational speed [rpm]};
    		\end{tikzpicture}
    	\caption{Example of a drawing + explanations created with TikZ, here: transfer of mechanical power with a shaft}\label{fig:torque}
    \end{figure}
\end{frame}

\section{Plotting airfoils}
The Applied Aerodynamics Group of the university of Illinois has an online database of more than 1600 airfoils \citep{airfoils}. The airfoil data is made available in data files with the extension \verb|.dat|. Data files can be plotted directly in \LaTeX\ if the headers are removed. See \cref{fig:airfoil} where the \href{https://m-selig.ae.illinois.edu/ads/coord/e625.dat}{Eppler 625 airfoil data} is plotted, combined with a grid.

\begin{frame}{Plotting of airfoil data contained in a file}
	\begin{figure}[htbp]\centering
			\begin{tikzpicture}[scale=0.9]
			\draw [help lines] (-1,-1) grid (11,2);
			\draw[scale=10,brown,thick] plot file{data/e625profile.dat} -- cycle;
			\node[brown,fill=white,right] at (0,1.5) {Eppler 625 airfoil};
			\end{tikzpicture}
		\caption{A plot of an Eppler 625 airfoil, to illustrate the plotting of data points from a file, data from \textcite{airfoils}}\label{fig:airfoil}
	\end{figure}
\end{frame}

\section{Plotting analytical functions}
It's also easy to plot functions in axes as shown in \cref{fig:analyticalfunctions}. In this example, three analytical functions are plotted. Inspect the source code to verify that these functions are analytically given.

\begin{figure}[htbp]\centering
    \def\funA{4*x*x-20*x-50} % definition of function A
    %\def\a{4}\def\b{-20}\def\c{-50}
        \begin{tikzpicture}
            \begin{axis}[
                    xmin=-6.28,xmax=6.28,
                    samples=200,
                    legend pos= outer north east,
                    legend cell align=left,
                    ]
                  \addplot[blue,ultra thin](x,{\funA}); % braces needed to protect
                  \addlegendentry{$y=4x^2-20x-50$};
                  \addplot[red,dashed](x,{100*sin(x*180/3.14)}); % braces needed to protect
                  \addlegendentry{$y=100\sin x$};
                  \addplot[olive, thick]({0.2*x*x*x},-40*x); % braces needed to protect
                  \addlegendentry{$x=f(y) \rightarrow$ see source code};
            \end{axis}
        \end{tikzpicture}
	\caption{Example of the plotting of some analytical functions}\label{fig:analyticalfunctions}
\end{figure}

\section{Plotting experimental data against a model}
It is also possible to plot data points together with an analytical model. This is shown in \cref{fig:psat} where experimental data is plotted against an parabolic model.

\begin{figure}[htbp]\centering
	\begin{tikzpicture}[scale=1.1]
		\begin{axis} [
		xlabel=$T\ \unitschart{\si{\celsius}}$,
		ylabel= $\psat\unitschart{kPa}$,
		ylabel near ticks,
		legend pos= north west,
		legend cell align=left,
		]
		\addplot[blue,domain=0:50,samples=100]{0.004*x^2+1};
        \addlegendentry{$\psat=0.004\cdot T^2+1$ (model)}
		\addplot[red,only marks, mark=x,error bars/y dir =both,error bars/y fixed = 1] coordinates {
			(0,0.61) (5,0.87) (10,1.23) (15,1.71) (20,2.34) (25,3.17) (30,4.25) (35,5.6286) (40,7.39) (45,9.5944) (50,12.352)};
			\addlegendentry{experiment};
		\end{axis}
	\end{tikzpicture}
	\caption{Example of a plot of an analytical model against experimental data with error bars; here: the saturation pressure of water as function of temperature}\label{fig:psat}
\end{figure}

\section{Predefining plot settings}
It's also possible to predefine the settings of charts. See for instance the definition of \verb|axespsychro| in the style file \verb|myDrawingStuff.sty| and how it is used as an option in: \verb|\begin{axis}[axespsychro]| in the source code of \cref{fig:psychrochart} and \cref{fig:psychrochart2}. Both charts have exactly the same appearance because of the same predefined settings.

\begin{figure}[ht]\centering
		\begin{tikzpicture}[scale=1]
    		\begin{axis}[axespsychro]
        		\drawsaturationline{sat};
        		\coordinate (A) at (axis cs: 30,10);
        		\coordinate (B) at (axis cs: 19.2,14.3);
        		\coordinate (C) at (axis cs: 13.5,10);
        		\draw [process,deco=0.6] (A) -- (B) node [above left] {$T_\up{WB}$};  
        		\draw [process,deco=0.6] (A) -- (C) node [above left] {$T_\up{DP}$}; 
        		\draw [helpline] (A) -- (axis cs:30,0) node [above right] {$T_\up{DB}$}; 
        		\draw [helpline] (A) -- (axis cs:50,10); % 
        		\drawdot{A};\drawdot{B};\drawdot{C};
    		\end{axis}
		\end{tikzpicture}
		\caption{Example of a simplified psychrometric chart with predefined settings; here: indication of the dew point and the wet bulb temperature}\label{fig:psychrochart}
\end{figure}

\begin{figure}[ht]\centering
		\begin{tikzpicture}[scale=1]
			\begin{axis}[axespsychro]
				\drawsaturationline{sat};
				\coordinate (A) at (axis cs: 40,5);
                \coordinate (B) at (axis cs: 4.5,5);
				\draw [process,deco=0.5] (A) -- (B) node [above left] {$T_\up{DP}$}; % 
				\draw [helpline] (A) -- (axis cs: 40,0); % 
				\draw [helpline] (A) -- (axis cs: 50,5); % 
				\drawdot{A};\drawdot{B};
				%\tekenisoHlijnenoppsychrokaart; % teken isenthalpen
			\end{axis}
		\end{tikzpicture}
		\caption{Example of a second psychrometric chart with the same predefined settings; here: indication of the dew point}\label{fig:psychrochart2}
\end{figure}

\section{Grouping of plots}
Plots can be grouped in such a way that they share the same axes, see example in \cref{fig:groupplot}.

\begin{frame}{An example of a group plot with shared axes}
	\begin{figure}[ht]\centering
		\input{figs/groupplot} % if the code for a drawing is long, you can opt to put it in a separate file
		\caption{The axes can be shared in group plots, here: impact of surface area on heat exchanger performance}\label{fig:groupplot}
	\end{figure}
\end{frame}

\section{Plotting data from an external file}
Data can also be read from external files and then plotted, as in \cref{fig:thermalpower}. Even if the data in the file contains date- and timestamp (like for instance of the type \texttt{YYYY-MM-DD HH:MM}) they can still be plotted by using the TikZ-library \texttt{dateplot}. An example of such a plot is shown in \cref{fig:weatherdata}. The meteo data was fetched from \href{https://mesonet.agron.iastate.edu/request/download.phtml?}{Iowa Environmental Mesonet (IEM)} from Iowa State University. The IEM maintains an ever growing archive of automated airport weather observations from around the world.

\begin{frame}{Example of plotting data from a file}
    \begin{figure}[ht]\centering
        \begin{tikzpicture}[scale=0.8]
          \begin{axis}[
              width=0.9\linewidth, % Scale the plot to \linewidth
              grid=major, 
              grid style={dashed,gray!30},
              xlabel={Time}, % Set the labels
              ylabel={Thermal power},
              x unit=\si{\hour}, % Set the respective units
              y unit=\si{\kilo\watt}, % Set the respective units
              legend pos= north east,
              ]
              \addplot[cyan,mark=x] table[x index=0, y index=1,col sep=comma] {data/data_heat.csv}; % read data from file
              \legend{building H}
          \end{axis}
        \end{tikzpicture}
        \caption{Example of plotted data read from an external csv-file; here: thermal power demand from a building during one week}\label{fig:thermalpower}
    \end{figure}
\end{frame}

\begin{frame}{Example of plotting data with dates from a file}
    \begin{figure}[ht]\centering
        \begin{tikzpicture}[scale=0.9]
              \begin{axis}[
                width=0.8\linewidth,
                date coordinates in=x, % to indicate that the dates are in x-axis
                %xmin=2020-03-01 00:00,
                %xmax=2020-03-07 23:59,
                xtick distance=1, % set `xtick distance' to 24 hours (24/24)
                % xtick=data, % to use every entry as xtick
                xticklabel={\day-\month},
		        xlabel={Temperature data for Brookneal Virginia (USA), first week of April '20},
                ylabel={Temperature},
                y unit=\si{\celsius}, % Set the respective units
                legend pos=south east, % to place the legend
                legend cell align=left, % to align the text in the legend
                ]
                \addplot [red] table [col sep=comma,x index=1, y index=2] {data/weatherVirginia.csv};
                \addplot [blue] table [col sep=comma,x index=1, y index=3] {data/weatherVirginia.csv};
                \legend{dry bulb temperature, dewpoint}
              \end{axis}
        \end{tikzpicture}
        \caption{Example of data with date- and timestamps plotted with the TikZ-library \texttt{dateplot}; here: some weather data during one week}\label{fig:weatherdata}
    \end{figure}
\end{frame}

\mode<all> 
% Required because of the \mode* command at the beginning of this file.
% See beameruserfile.pdf section 21.3 for more info